We develop a rigorous algebraic framework for chain-of-thought (CoT) prompt
composition, modeled as a semigroup of linear operators on a metric space of
reasoning trajectories. A CoT operator $T(A,B)$ acts on a trajectory
$\tau=(s_0,s_1,\dots,s_n)$ by preserving the initial query~$s_0$ and
transforming subsequent states via $s_i\mapsto As_{i-1}+Bs_i$, where $A$
encodes context propagation and $B$ encodes current-state refinement. We prove
five main results: (1)~the set of CoT operators is closed under composition
and forms a monoid with an explicit composition formula; (2)~a sharp
contractivity criterion showing that iterates $T^n$ converge in the strong
operator topology if and only if the spectral radius $\rho(B)<1$, with
convergence rate $O(\rho(B)^n)$; (3)~idempotent CoT operators correspond
exactly to projection operators satisfying $B^2=B$ and $BA=0$; (4)~a complete
four-case classification of convergence behavior---contractive, projective,
oscillatory, and divergent---determined by the spectrum of~$B$; and (5)~the
Jacobs--de~Leeuw--Glicksberg decomposition separates reasoning trajectories
into stable reasoning patterns and transient noise. All results are verified
computationally on explicit finite-dimensional examples. The framework
transforms CoT prompting from a purely empirical technique into a
mathematically principled methodology with actionable convergence criteria.
